% 検証（クラス幅）：JFMクラス幅の実装（約物・半角の字送り縮小）
% class@<n>{w}{h}{d}{align} の w<1zw クラスは \hbox to w*zw で字送りを
% 縮小し，align に従ってグリフを配置する（ltj の align 相当）．
% 測定: typeout で実測幅を出力し，pairs.tex 形式の PASS/FAIL 判定も行う．
\documentclass[pdflatex]{jlreq}
\begin{document}
\makeatletter
\newbox\pwbox
\def\pwmeas#1#2#3{%
  \ifx d\jlreq@engine\global\jlreqpdf@lastclass=90\fi
  \setbox\pwbox=\hbox{#3}%
  \ifdim\wd\pwbox=#1\relax
    \typeout{PASS: #2 = \the\wd\pwbox}%
  \else
    \typeout{FAIL: #2 expected #1, got \the\wd\pwbox}%
  \fi
}
% 単字: 約物は 0.5zw，かな/漢字/記号類は 1zw
\pwmeas{5.0pt}{touten-halfwidth (class7 w=0.5 align=left)}{、}
\pwmeas{5.0pt}{kuten-halfwidth (class6 w=0.5 align=left)}{。}
\pwmeas{5.0pt}{close-bracket-halfwidth (class2 w=0.5 align=left)}{」}
\pwmeas{5.0pt}{colon-halfwidth (class5 w=0.5 align=middle)}{：}
% 90->1 = {0,0,0}（parbdd突き出し: インクは左へはみ出し，字送りは0.5zw）
\pwmeas{5.0pt}{open-bracket-tentsuki (class1 w=0.5 align=right, 90->1=0)}{「}
\pwmeas{10.0pt}{kanji-fullwidth (class0 w=1)}{漢}
\pwmeas{10.0pt}{prolonged-fullwidth (class10 w=1 align=middle)}{ー}
\pwmeas{10.0pt}{excl-fullwidth (class4 w=1 align=middle)}{！}
% 組合せ: ペアグルーはグリフ箱と独立
% 15->7 = nat 0 (kernmode抑制) : 10+0+5
\pwmeas{15.0pt}{kanji-touten (15->7 nat 0)}{あ、}
% 7->15 = nat 0.5zw : 5+5+10
\pwmeas{20.0pt}{touten-kanji (7->15 nat 0.5zw)}{、あ}
% 15->1 = nat 0.5zw + 開き括弧右寄せ : 10+5+5
\pwmeas{20.0pt}{kanji-openbracket (15->1 nat 0.5zw)}{あ「}
% 2->15 = nat 0.5zw : 5+5+10
\pwmeas{20.0pt}{closebracket-kanji (2->15 nat 0.5zw)}{」あ}
\pwmeas{20.0pt}{kanji-kanji (kanjiskip nat 0)}{漢字}
\makeatother
% 実ページ（目視確認用）: 約物の詰まりと開き括弧の突き出し
\parindent=0pt
あ、い。う「え」お：か！き…くー
「x」（y）
\end{document}
